% 第一章 绪言
\chapter{绪言}

\section{研究背景}

随着移动互联网、本地生活服务平台与即时配送体系的成熟，用户在餐饮选择上拥有了前所未有的丰富供给。以美团、饿了么、大众点评为代表的本地生活平台聚合了数以万计的餐饮商户，用户在手机端即可浏览海量菜品信息。对于高校场景而言，食堂、校内商户、校外餐馆、外卖平台与便利食品共同构成了高密度、多样化的候选集合。表面上看，候选数量增加意味着用户有更大概率找到符合口味与价格预期的餐品；但从决策行为角度分析，候选扩张同时也显著提高了用户的选择成本，尤其在时间紧迫、预算有限或身体状态波动明显时，这种成本会进一步放大。

诺贝尔经济学奖得主赫伯特·西蒙（Herbert Simon）提出的"有限理性"理论指出，决策者在信息过载环境中难以对所有选项进行完全理性的评估，往往会采用启发式策略或陷入决策瘫痪\textsuperscript{[1]}。在餐饮选择场景中，这一现象表现得尤为明显：当用户面对数十甚至上百个候选菜品时，"选择困难"不仅消耗时间，还会带来心理负担。特别是在高校场景中，学生群体具有用餐时间集中、预算敏感、健康意识增强等特点，对推荐系统的即时性和可执行性提出了更高要求。

餐饮推荐与传统商品推荐存在明显差异。商品推荐（如电商平台的商品推荐）通常以长期偏好和历史行为为主要依据，用户购买决策周期较长，且对即时约束的要求相对较低。而餐饮推荐更强调即时决策与现实约束的协调：用户在同一天内可能因为时段、天气、情绪、身体状态或同伴关系而改变饮食偏好。例如，早餐倾向清淡，夜宵偏好高热量或重口味；减脂阶段更关注低油低糖，胃部不适时则倾向温和饮食；单人用餐追求便捷，聚餐场景则注重菜品的丰富性与分享性。因此，单纯依靠历史点击或购买记录进行推荐，往往会出现"偏好相近但场景不适"的问题，导致推荐结果在理论上合理、在现实中却难以执行。

在此背景下，构建一套能够显式捕捉当前场景、并对预算、禁忌与距离进行约束优化的餐饮推荐系统，具有较强的现实意义。本文所研究的问题并非泛化意义上的"美食推荐"，而是面向"当前想吃什么"这一即时决策任务，目标是在较短交互成本下输出可执行、可解释且场景一致的推荐结果。

\section{国内外研究现状}

\subsection{推荐系统基础方法研究}
经过多年发展，推荐系统已形成较为成熟的技术体系。在早期研究中，协同过滤与基于内容的推荐是两种最具代表性的方法。协同过滤依托用户与物品之间的交互行为挖掘相似关系，进而完成个性化推荐，其中又分为基于用户的协同过滤与基于物品的协同过滤两类典型实现方式\super{[2]}。不过，协同过滤方法高度依赖历史交互数据，在用户行为稀疏或新用户、新物品较多的场景下，推荐效果会出现明显下降。基于内容的推荐则以物品自身属性为依据，通过提取特征并匹配用户历史偏好完成推荐，可解释性较强，但也容易陷入推荐范围固化、结果同质化的问题，难以发掘用户的潜在兴趣\super{[3]}。

随着数据规模扩大与深度学习技术的普及，矩阵分解、神经协同过滤、注意力机制以及图神经网络等方法被广泛引入推荐领域，进一步提升了模型的特征表达与泛化能力。Koren等人提出的矩阵分解技术，通过对用户—物品交互矩阵进行低维分解，有效改善了数据稀疏带来的影响\super{[4]}。He等人提出的神经协同过滤模型，使用多层感知机替代传统线性内积，让模型能够更好地拟合用户与物品之间的非线性交互关系\super{[5]}。近几年，基于Transformer的序列推荐模型得到快速发展，例如SASRec、BERT4Rec等方法借助自注意力机制捕捉用户行为序列中的时序依赖，在公开数据集上展现出更优的推荐性能\super{[6]}。

\subsection{餐饮推荐与本地生活推荐研究}
在餐饮及本地生活服务场景中，推荐思路逐渐从静态偏好建模转向上下文感知推荐。与传统推荐场景不同，餐饮消费受时间、位置、环境、社交关系等实时因素影响明显，将这些情境信息融入模型，能够有效提升推荐结果的实用性。Adomavicius等人将上下文信息归纳为用户、物品与交互三大类别，并搭建了多维上下文感知的推荐框架\super{[7]}。

在地理位置相关研究中，用户与商家的空间距离是影响餐饮选择的重要因素。Liu等人在矩阵分解模型中加入地理影响因子，利用距离衰减规律刻画用户的实际活动范围\super{[8]}。Yin等人基于用户历史签到数据建模位置偏好，证实了餐饮选择中存在显著的空间聚集特征\super{[9]}。在文本特征利用方面，菜品名称、口味标签、食材构成与描述文本都可作为有效特征。Zhang等人采用Word2Vec对菜品文本进行向量表示，结合用户评论文本构建语义匹配模型，使召回结果更贴合用户偏好\super{[10]}。在健康导向推荐方面，部分研究将热量、营养成分等指标加入优化目标，为有特殊饮食需求的用户提供更具针对性的推荐服务\super{[11]}。

\subsection{交互式推荐与问答引导研究}
交互式推荐与对话式推荐是近年来的研究热点，其核心思路是通过主动交互补充用户偏好信息，而不是被动依赖历史行为。在用户信息不足时，系统可通过提问、对话等方式快速获取关键偏好，快速构建有效的推荐依据。Christakopoulou等人基于强化学习设计对话推荐策略，动态选择提问内容以最大化信息获取效率\super{[12]}。Zhang等人提出基于属性选择的交互推荐框架，通过多轮提问逐步缩小候选范围，提升推荐精准度\super{[13]}。

在餐饮推荐场景中，问答式交互具备天然优势。就餐相关的约束条件相对明确，包括预算、口味、距离、饮食禁忌等，用户通常可以快速完成结构化作答。相关研究表明，仅通过3—5个关键问题的轻量交互，就能明显缩小候选范围并提升用户体验\super{[14]}。但目前多数交互式餐饮推荐工作仍停留在算法验证阶段，对前后端系统联动、交互日志记录与工程落地等实际问题考虑不足，距离实际可用的系统仍有差距。

\subsection{现有研究的不足与本文定位}

综合上述研究现状，现有工作存在以下不足：

\begin{enumerate}[label=\arabic*., leftmargin=2em, itemsep=0.2em]
    \item \textbf{约束建模与偏好建模脱节}：多数系统侧重于提升推荐结果的偏好匹配度，而对预算、距离、禁忌等硬性约束的执行性关注不足，导致推荐结果"好看但不可用"。
    \item \textbf{长期偏好与即时场景割裂}：传统推荐系统过度依赖历史行为，难以适应用户在同一天内的偏好波动；而纯问答系统又忽略了用户的长期特征，导致推荐缺乏个性化基础。
    \item \textbf{工程闭环不完整}：很多研究停留在算法原型阶段，缺少完整的前端交互、后端服务、数据存储与日志回流机制，难以形成可持续迭代的产品。
\end{enumerate}

针对上述问题，本文提出了一种融合静态画像与即时问答的双通道建模方法，并采用"硬性过滤---向量召回---综合重排"的三阶段推荐策略，在保证约束可执行的前提下提升语义相关性。同时，本文完成了微信小程序前端、FastAPI 后端、MySQL 数据存储与向量化模块的完整工程实现，形成了从需求采集到结果反馈的闭环链路。

\section{本文的研究目标和内容}

为了让高校及周边用户在面对"今天吃什么"这一高频决策问题时能够快速获得符合当前场景的可执行推荐，本文设计并实现了一套基于用户画像与问答引导的智能推荐微信小程序。该系统包括面向用户的微信推荐小程序、基于 FastAPI 的后端服务与数据管理平台，以及融合语义向量与约束排序的智能推荐引擎。

本文主要研究内容包括以下几个方面：

\textbf{（1）设计并实现微信推荐小程序。}通过分析高校用户在餐饮选择中的核心痛点（候选过载、场景多变、约束复杂），确定小程序以"即时场景采集+可执行推荐"为核心定位；学习微信小程序开发技术，设计"首页---问答---推荐---确认---历史"的闭环交互流程，实现用户注册登录、即时问答引导、推荐结果展示、选择确认与历史回看等功能；结合微信定位 API，支持基于地理位置的附近美食浏览，为用户提供从决策到执行的完整链路。

\textbf{（2）设计并实现后端服务与数据管理平台。}基于 FastAPI 框架构建高性能的 RESTful API 服务，按业务职责划分为认证、画像、菜品、推荐与日志五大模块；使用 SQLAlchemy ORM 完成数据库建模，采用 MySQL 持久化存储用户、画像、菜品、店铺、历史记录和交互日志；设计基于 JWT 的鉴权机制保障接口安全，实现前后端分离的分层架构，使系统具备良好的可维护性与可扩展性。

\textbf{（3）研究基于用户画像与问答引导的推荐方法，分析传统推荐算法在即时餐饮场景中的不足。}在研究语义向量召回与约束排序的基础上，针对传统协同过滤和基于历史行为的推荐方法在"冷启动""场景失配"等方面的问题，提出"静态画像+即时问答"的双通道建模思路；设计"硬性过滤---向量召回---综合重排"的三段式推荐策略，将预算、禁忌、距离等硬性约束前置执行，再通过语义向量捕捉用户当前意图，最后融合距离惩罚与历史平滑进行综合排序，使推荐结果兼顾可执行性与场景一致性。

\textbf{（4）设计并实现菜品向量化与语义召回模块。}在采集菜品的多维属性信息（菜名、菜系、口味标签、食材、描述等）后，基于 sentence-transformers 本地模型将结构化字段拼接为语义文本并编码为稠密向量；融合用户静态画像（长期偏好、健康目标、饮食禁忌）与即时问答（用餐时段、预算、口味强度、特殊状态）构建用户目标向量，通过余弦相似度计算实现语义召回；引入特殊状态加权机制（如减脂、胃不舒服等场景），对召回结果进行场景化增强，提升推荐结果与当前状态的匹配度。

\textbf{（5）设计并实现数据闭环与测试验证机制。}建立交互日志体系，记录推荐曝光、点击、确认和负样本行为，形成"采集---推荐---反馈---回流"的数据闭环；开展功能测试验证各模块正确性，开展场景化测试验证推荐算法在不同预算、时段、特殊状态下的适配能力，开展性能测试与鲁棒性测试验证系统的响应速度与异常处理能力，确保系统具备工程可用性。

\section{论文组织结构}

全文共五章，各章内容安排如下：

\textbf{第一章 绪言}。介绍研究背景、国内外研究现状、研究意义与目的、研究内容与技术路线，以及论文组织结构。

\textbf{第二章 系统相关技术概述}。介绍推荐系统基本原理、微信小程序开发技术、FastAPI 后端框架、SQLAlchemy ORM 技术、MySQL 数据库、sentence-transformers 向量化技术以及 JWT 鉴权机制，并对关键技术选型进行对比分析。

\textbf{第三章 系统设计与实现}。阐述系统需求分析、总体架构设计、数据库设计、核心模块详细设计（认证、画像、菜品、推荐、历史日志）、前端交互实现，并给出关键流程图与界面说明。

\textbf{第四章 系统功能测试}。给出测试目标与测试环境、系统功能性测试、场景化推荐测试、非功能性测试（性能、可靠性、易用性、鲁棒性）以及测试结论。

\textbf{第五章 总结与展望}。总结全文工作，分析主要特点与创新点，指出存在的问题，并提出后续改进方向。
